\documentclass[11pt,a4paper]{article}

% shared/macros.tex
% Shared notation for all surface-partition math documents.
% Include with: % shared/macros.tex
% Shared notation for all surface-partition math documents.
% Include with: % shared/macros.tex
% Shared notation for all surface-partition math documents.
% Include with: \input{../shared/macros}

% -------------------------------------------------------------------------
% Packages (include here so each document only needs \input{../shared/macros})
% -------------------------------------------------------------------------
\usepackage{amsmath,amssymb,amsthm}
\usepackage{bm}           % \bm for bold math
\usepackage{mathtools}    % \coloneqq, dcases, etc.
\usepackage{booktabs}     % \toprule etc. in tables
\usepackage{hyperref}
\usepackage{geometry}
\usepackage{microtype}
\usepackage{parskip}
\usepackage{xcolor}
\usepackage{tcolorbox}
\usepackage{enumitem}
\usepackage{float}             % [H] float placement

\geometry{a4paper, margin=2.5cm}

% -------------------------------------------------------------------------
% Theorem-like environments
% -------------------------------------------------------------------------
\theoremstyle{definition}
\newtheorem{defn}{Definition}[section]
\newtheorem{remark}{Remark}[section]

\theoremstyle{plain}
\newtheorem{prop}{Proposition}[section]

% -------------------------------------------------------------------------
% Mesh and geometry
% -------------------------------------------------------------------------
\newcommand{\V}{\mathbf{V}}            % vertex array V[i] in R^dim
\newcommand{\F}{\mathbf{F}}            % face (triangle) array
\newcommand{\vone}[1]{v_{1,#1}}       % first vertex index of VP i
\newcommand{\vtwo}[1]{v_{2,#1}}       % second vertex index of VP i

% -------------------------------------------------------------------------
% Variable points
% -------------------------------------------------------------------------
\newcommand{\lam}{\lambda}             % scalar VP parameter
\newcommand{\Lam}{\boldsymbol{\lambda}}% full parameter vector
\newcommand{\vp}[1]{\mathbf{p}_{#1}}  % 3D position of VP i
\newcommand{\ed}[1]{\mathbf{d}_{#1}}  % edge direction d_i = V[v1_i] - V[v2_i]

% -------------------------------------------------------------------------
% Steiner / triple-point
% -------------------------------------------------------------------------
\newcommand{\Sv}{\mathbf{S}}           % Steiner (Fermat-Torricelli) point
\newcommand{\Ki}{K_{i}}               % projection matrix (I-n_i n_i^T)/r_i
\newcommand{\Mmat}{M}                 % sum of K_i matrices at a triple point

% -------------------------------------------------------------------------
% Normals and unit vectors
% -------------------------------------------------------------------------
\newcommand{\nhat}{\hat{\mathbf{n}}}  % unit normal to a triangle
\newcommand{\Dhat}{\hat{\boldsymbol{\Delta}}}  % unit segment direction

% -------------------------------------------------------------------------
% Operations
% -------------------------------------------------------------------------
\newcommand{\norm}[1]{\left\|#1\right\|}
\newcommand{\abs}[1]{\left|#1\right|}
\newcommand{\ip}[2]{\langle #1,\, #2 \rangle}

% Projection perpendicular to a unit vector \uhat
\newcommand{\Proj}[1]{\bigl(\mathbf{I} - #1 #1^\top\bigr)}

% argmax operator (winner-take-all assignment in Phase 1 documents)
\DeclareMathOperator*{\argmax}{arg\,max}
\DeclareMathOperator*{\argmin}{arg\,min}

% -------------------------------------------------------------------------
% Calculus shorthands
% -------------------------------------------------------------------------
\newcommand{\pd}[2]{\dfrac{\partial #1}{\partial #2}}
\newcommand{\pdd}[3]{\dfrac{\partial^2 #1}{\partial #2\,\partial #3}}
\newcommand{\grad}{\nabla}
\newcommand{\Hess}[1]{\nabla^2 #1}

% -------------------------------------------------------------------------
% Optimization problem objects
% -------------------------------------------------------------------------
\newcommand{\Preg}{P_{\mathrm{reg}}}  % regular perimeter
\newcommand{\Pst}{P_{\mathrm{st}}}   % Steiner perimeter
\newcommand{\Areg}{A^{\mathrm{reg}}} % regular cell area
\newcommand{\Ast}{A^{\mathrm{st}}}   % Steiner area contribution
\newcommand{\Atgt}{\bar{A}}           % target area per cell
\newcommand{\Lagr}{\mathcal{L}}       % Lagrangian
\newcommand{\objfac}{\sigma}          % IPOPT objective scaling factor

% -------------------------------------------------------------------------
% Code cross-references (for "In the code..." callouts)
% -------------------------------------------------------------------------
\newcommand{\code}[1]{\texttt{#1}}
\newcommand{\codefile}[1]{\texttt{\small #1}}

% -------------------------------------------------------------------------
% Threshold dynamics / approach B (10-mbo-auction-dynamics)
% -------------------------------------------------------------------------
\newcommand{\Dlump}{D}                 % lumped mass matrix diag(v)
\newcommand{\Atau}{A_\tau}             % discrete diffusion operator (M + tau K)^{-1} M
\newcommand{\Etau}{E_\tau}             % heat-content (threshold-dynamics) energy
\newcommand{\Rcell}{R_{\mathrm{cell}}} % radius of the equal-area disc sqrt(Abar/pi)
\newcommand{\hmean}{h_{\mathrm{mean}}} % mean edge length
\newcommand{\hmax}{h_{\mathrm{max}}}   % max edge length
\newcommand{\Gt}{G_t}                  % heat kernel at time t
\newcommand{\Ktau}{K^{\mathrm{res}}_\tau} % resolvent (backward-Euler) kernel


% -------------------------------------------------------------------------
% Packages (include here so each document only needs % shared/macros.tex
% Shared notation for all surface-partition math documents.
% Include with: \input{../shared/macros}

% -------------------------------------------------------------------------
% Packages (include here so each document only needs \input{../shared/macros})
% -------------------------------------------------------------------------
\usepackage{amsmath,amssymb,amsthm}
\usepackage{bm}           % \bm for bold math
\usepackage{mathtools}    % \coloneqq, dcases, etc.
\usepackage{booktabs}     % \toprule etc. in tables
\usepackage{hyperref}
\usepackage{geometry}
\usepackage{microtype}
\usepackage{parskip}
\usepackage{xcolor}
\usepackage{tcolorbox}
\usepackage{enumitem}
\usepackage{float}             % [H] float placement

\geometry{a4paper, margin=2.5cm}

% -------------------------------------------------------------------------
% Theorem-like environments
% -------------------------------------------------------------------------
\theoremstyle{definition}
\newtheorem{defn}{Definition}[section]
\newtheorem{remark}{Remark}[section]

\theoremstyle{plain}
\newtheorem{prop}{Proposition}[section]

% -------------------------------------------------------------------------
% Mesh and geometry
% -------------------------------------------------------------------------
\newcommand{\V}{\mathbf{V}}            % vertex array V[i] in R^dim
\newcommand{\F}{\mathbf{F}}            % face (triangle) array
\newcommand{\vone}[1]{v_{1,#1}}       % first vertex index of VP i
\newcommand{\vtwo}[1]{v_{2,#1}}       % second vertex index of VP i

% -------------------------------------------------------------------------
% Variable points
% -------------------------------------------------------------------------
\newcommand{\lam}{\lambda}             % scalar VP parameter
\newcommand{\Lam}{\boldsymbol{\lambda}}% full parameter vector
\newcommand{\vp}[1]{\mathbf{p}_{#1}}  % 3D position of VP i
\newcommand{\ed}[1]{\mathbf{d}_{#1}}  % edge direction d_i = V[v1_i] - V[v2_i]

% -------------------------------------------------------------------------
% Steiner / triple-point
% -------------------------------------------------------------------------
\newcommand{\Sv}{\mathbf{S}}           % Steiner (Fermat-Torricelli) point
\newcommand{\Ki}{K_{i}}               % projection matrix (I-n_i n_i^T)/r_i
\newcommand{\Mmat}{M}                 % sum of K_i matrices at a triple point

% -------------------------------------------------------------------------
% Normals and unit vectors
% -------------------------------------------------------------------------
\newcommand{\nhat}{\hat{\mathbf{n}}}  % unit normal to a triangle
\newcommand{\Dhat}{\hat{\boldsymbol{\Delta}}}  % unit segment direction

% -------------------------------------------------------------------------
% Operations
% -------------------------------------------------------------------------
\newcommand{\norm}[1]{\left\|#1\right\|}
\newcommand{\abs}[1]{\left|#1\right|}
\newcommand{\ip}[2]{\langle #1,\, #2 \rangle}

% Projection perpendicular to a unit vector \uhat
\newcommand{\Proj}[1]{\bigl(\mathbf{I} - #1 #1^\top\bigr)}

% argmax operator (winner-take-all assignment in Phase 1 documents)
\DeclareMathOperator*{\argmax}{arg\,max}
\DeclareMathOperator*{\argmin}{arg\,min}

% -------------------------------------------------------------------------
% Calculus shorthands
% -------------------------------------------------------------------------
\newcommand{\pd}[2]{\dfrac{\partial #1}{\partial #2}}
\newcommand{\pdd}[3]{\dfrac{\partial^2 #1}{\partial #2\,\partial #3}}
\newcommand{\grad}{\nabla}
\newcommand{\Hess}[1]{\nabla^2 #1}

% -------------------------------------------------------------------------
% Optimization problem objects
% -------------------------------------------------------------------------
\newcommand{\Preg}{P_{\mathrm{reg}}}  % regular perimeter
\newcommand{\Pst}{P_{\mathrm{st}}}   % Steiner perimeter
\newcommand{\Areg}{A^{\mathrm{reg}}} % regular cell area
\newcommand{\Ast}{A^{\mathrm{st}}}   % Steiner area contribution
\newcommand{\Atgt}{\bar{A}}           % target area per cell
\newcommand{\Lagr}{\mathcal{L}}       % Lagrangian
\newcommand{\objfac}{\sigma}          % IPOPT objective scaling factor

% -------------------------------------------------------------------------
% Code cross-references (for "In the code..." callouts)
% -------------------------------------------------------------------------
\newcommand{\code}[1]{\texttt{#1}}
\newcommand{\codefile}[1]{\texttt{\small #1}}

% -------------------------------------------------------------------------
% Threshold dynamics / approach B (10-mbo-auction-dynamics)
% -------------------------------------------------------------------------
\newcommand{\Dlump}{D}                 % lumped mass matrix diag(v)
\newcommand{\Atau}{A_\tau}             % discrete diffusion operator (M + tau K)^{-1} M
\newcommand{\Etau}{E_\tau}             % heat-content (threshold-dynamics) energy
\newcommand{\Rcell}{R_{\mathrm{cell}}} % radius of the equal-area disc sqrt(Abar/pi)
\newcommand{\hmean}{h_{\mathrm{mean}}} % mean edge length
\newcommand{\hmax}{h_{\mathrm{max}}}   % max edge length
\newcommand{\Gt}{G_t}                  % heat kernel at time t
\newcommand{\Ktau}{K^{\mathrm{res}}_\tau} % resolvent (backward-Euler) kernel
)
% -------------------------------------------------------------------------
\usepackage{amsmath,amssymb,amsthm}
\usepackage{bm}           % \bm for bold math
\usepackage{mathtools}    % \coloneqq, dcases, etc.
\usepackage{booktabs}     % \toprule etc. in tables
\usepackage{hyperref}
\usepackage{geometry}
\usepackage{microtype}
\usepackage{parskip}
\usepackage{xcolor}
\usepackage{tcolorbox}
\usepackage{enumitem}
\usepackage{float}             % [H] float placement

\geometry{a4paper, margin=2.5cm}

% -------------------------------------------------------------------------
% Theorem-like environments
% -------------------------------------------------------------------------
\theoremstyle{definition}
\newtheorem{defn}{Definition}[section]
\newtheorem{remark}{Remark}[section]

\theoremstyle{plain}
\newtheorem{prop}{Proposition}[section]

% -------------------------------------------------------------------------
% Mesh and geometry
% -------------------------------------------------------------------------
\newcommand{\V}{\mathbf{V}}            % vertex array V[i] in R^dim
\newcommand{\F}{\mathbf{F}}            % face (triangle) array
\newcommand{\vone}[1]{v_{1,#1}}       % first vertex index of VP i
\newcommand{\vtwo}[1]{v_{2,#1}}       % second vertex index of VP i

% -------------------------------------------------------------------------
% Variable points
% -------------------------------------------------------------------------
\newcommand{\lam}{\lambda}             % scalar VP parameter
\newcommand{\Lam}{\boldsymbol{\lambda}}% full parameter vector
\newcommand{\vp}[1]{\mathbf{p}_{#1}}  % 3D position of VP i
\newcommand{\ed}[1]{\mathbf{d}_{#1}}  % edge direction d_i = V[v1_i] - V[v2_i]

% -------------------------------------------------------------------------
% Steiner / triple-point
% -------------------------------------------------------------------------
\newcommand{\Sv}{\mathbf{S}}           % Steiner (Fermat-Torricelli) point
\newcommand{\Ki}{K_{i}}               % projection matrix (I-n_i n_i^T)/r_i
\newcommand{\Mmat}{M}                 % sum of K_i matrices at a triple point

% -------------------------------------------------------------------------
% Normals and unit vectors
% -------------------------------------------------------------------------
\newcommand{\nhat}{\hat{\mathbf{n}}}  % unit normal to a triangle
\newcommand{\Dhat}{\hat{\boldsymbol{\Delta}}}  % unit segment direction

% -------------------------------------------------------------------------
% Operations
% -------------------------------------------------------------------------
\newcommand{\norm}[1]{\left\|#1\right\|}
\newcommand{\abs}[1]{\left|#1\right|}
\newcommand{\ip}[2]{\langle #1,\, #2 \rangle}

% Projection perpendicular to a unit vector \uhat
\newcommand{\Proj}[1]{\bigl(\mathbf{I} - #1 #1^\top\bigr)}

% argmax operator (winner-take-all assignment in Phase 1 documents)
\DeclareMathOperator*{\argmax}{arg\,max}
\DeclareMathOperator*{\argmin}{arg\,min}

% -------------------------------------------------------------------------
% Calculus shorthands
% -------------------------------------------------------------------------
\newcommand{\pd}[2]{\dfrac{\partial #1}{\partial #2}}
\newcommand{\pdd}[3]{\dfrac{\partial^2 #1}{\partial #2\,\partial #3}}
\newcommand{\grad}{\nabla}
\newcommand{\Hess}[1]{\nabla^2 #1}

% -------------------------------------------------------------------------
% Optimization problem objects
% -------------------------------------------------------------------------
\newcommand{\Preg}{P_{\mathrm{reg}}}  % regular perimeter
\newcommand{\Pst}{P_{\mathrm{st}}}   % Steiner perimeter
\newcommand{\Areg}{A^{\mathrm{reg}}} % regular cell area
\newcommand{\Ast}{A^{\mathrm{st}}}   % Steiner area contribution
\newcommand{\Atgt}{\bar{A}}           % target area per cell
\newcommand{\Lagr}{\mathcal{L}}       % Lagrangian
\newcommand{\objfac}{\sigma}          % IPOPT objective scaling factor

% -------------------------------------------------------------------------
% Code cross-references (for "In the code..." callouts)
% -------------------------------------------------------------------------
\newcommand{\code}[1]{\texttt{#1}}
\newcommand{\codefile}[1]{\texttt{\small #1}}

% -------------------------------------------------------------------------
% Threshold dynamics / approach B (10-mbo-auction-dynamics)
% -------------------------------------------------------------------------
\newcommand{\Dlump}{D}                 % lumped mass matrix diag(v)
\newcommand{\Atau}{A_\tau}             % discrete diffusion operator (M + tau K)^{-1} M
\newcommand{\Etau}{E_\tau}             % heat-content (threshold-dynamics) energy
\newcommand{\Rcell}{R_{\mathrm{cell}}} % radius of the equal-area disc sqrt(Abar/pi)
\newcommand{\hmean}{h_{\mathrm{mean}}} % mean edge length
\newcommand{\hmax}{h_{\mathrm{max}}}   % max edge length
\newcommand{\Gt}{G_t}                  % heat kernel at time t
\newcommand{\Ktau}{K^{\mathrm{res}}_\tau} % resolvent (backward-Euler) kernel


% -------------------------------------------------------------------------
% Packages (include here so each document only needs % shared/macros.tex
% Shared notation for all surface-partition math documents.
% Include with: % shared/macros.tex
% Shared notation for all surface-partition math documents.
% Include with: \input{../shared/macros}

% -------------------------------------------------------------------------
% Packages (include here so each document only needs \input{../shared/macros})
% -------------------------------------------------------------------------
\usepackage{amsmath,amssymb,amsthm}
\usepackage{bm}           % \bm for bold math
\usepackage{mathtools}    % \coloneqq, dcases, etc.
\usepackage{booktabs}     % \toprule etc. in tables
\usepackage{hyperref}
\usepackage{geometry}
\usepackage{microtype}
\usepackage{parskip}
\usepackage{xcolor}
\usepackage{tcolorbox}
\usepackage{enumitem}
\usepackage{float}             % [H] float placement

\geometry{a4paper, margin=2.5cm}

% -------------------------------------------------------------------------
% Theorem-like environments
% -------------------------------------------------------------------------
\theoremstyle{definition}
\newtheorem{defn}{Definition}[section]
\newtheorem{remark}{Remark}[section]

\theoremstyle{plain}
\newtheorem{prop}{Proposition}[section]

% -------------------------------------------------------------------------
% Mesh and geometry
% -------------------------------------------------------------------------
\newcommand{\V}{\mathbf{V}}            % vertex array V[i] in R^dim
\newcommand{\F}{\mathbf{F}}            % face (triangle) array
\newcommand{\vone}[1]{v_{1,#1}}       % first vertex index of VP i
\newcommand{\vtwo}[1]{v_{2,#1}}       % second vertex index of VP i

% -------------------------------------------------------------------------
% Variable points
% -------------------------------------------------------------------------
\newcommand{\lam}{\lambda}             % scalar VP parameter
\newcommand{\Lam}{\boldsymbol{\lambda}}% full parameter vector
\newcommand{\vp}[1]{\mathbf{p}_{#1}}  % 3D position of VP i
\newcommand{\ed}[1]{\mathbf{d}_{#1}}  % edge direction d_i = V[v1_i] - V[v2_i]

% -------------------------------------------------------------------------
% Steiner / triple-point
% -------------------------------------------------------------------------
\newcommand{\Sv}{\mathbf{S}}           % Steiner (Fermat-Torricelli) point
\newcommand{\Ki}{K_{i}}               % projection matrix (I-n_i n_i^T)/r_i
\newcommand{\Mmat}{M}                 % sum of K_i matrices at a triple point

% -------------------------------------------------------------------------
% Normals and unit vectors
% -------------------------------------------------------------------------
\newcommand{\nhat}{\hat{\mathbf{n}}}  % unit normal to a triangle
\newcommand{\Dhat}{\hat{\boldsymbol{\Delta}}}  % unit segment direction

% -------------------------------------------------------------------------
% Operations
% -------------------------------------------------------------------------
\newcommand{\norm}[1]{\left\|#1\right\|}
\newcommand{\abs}[1]{\left|#1\right|}
\newcommand{\ip}[2]{\langle #1,\, #2 \rangle}

% Projection perpendicular to a unit vector \uhat
\newcommand{\Proj}[1]{\bigl(\mathbf{I} - #1 #1^\top\bigr)}

% argmax operator (winner-take-all assignment in Phase 1 documents)
\DeclareMathOperator*{\argmax}{arg\,max}
\DeclareMathOperator*{\argmin}{arg\,min}

% -------------------------------------------------------------------------
% Calculus shorthands
% -------------------------------------------------------------------------
\newcommand{\pd}[2]{\dfrac{\partial #1}{\partial #2}}
\newcommand{\pdd}[3]{\dfrac{\partial^2 #1}{\partial #2\,\partial #3}}
\newcommand{\grad}{\nabla}
\newcommand{\Hess}[1]{\nabla^2 #1}

% -------------------------------------------------------------------------
% Optimization problem objects
% -------------------------------------------------------------------------
\newcommand{\Preg}{P_{\mathrm{reg}}}  % regular perimeter
\newcommand{\Pst}{P_{\mathrm{st}}}   % Steiner perimeter
\newcommand{\Areg}{A^{\mathrm{reg}}} % regular cell area
\newcommand{\Ast}{A^{\mathrm{st}}}   % Steiner area contribution
\newcommand{\Atgt}{\bar{A}}           % target area per cell
\newcommand{\Lagr}{\mathcal{L}}       % Lagrangian
\newcommand{\objfac}{\sigma}          % IPOPT objective scaling factor

% -------------------------------------------------------------------------
% Code cross-references (for "In the code..." callouts)
% -------------------------------------------------------------------------
\newcommand{\code}[1]{\texttt{#1}}
\newcommand{\codefile}[1]{\texttt{\small #1}}

% -------------------------------------------------------------------------
% Threshold dynamics / approach B (10-mbo-auction-dynamics)
% -------------------------------------------------------------------------
\newcommand{\Dlump}{D}                 % lumped mass matrix diag(v)
\newcommand{\Atau}{A_\tau}             % discrete diffusion operator (M + tau K)^{-1} M
\newcommand{\Etau}{E_\tau}             % heat-content (threshold-dynamics) energy
\newcommand{\Rcell}{R_{\mathrm{cell}}} % radius of the equal-area disc sqrt(Abar/pi)
\newcommand{\hmean}{h_{\mathrm{mean}}} % mean edge length
\newcommand{\hmax}{h_{\mathrm{max}}}   % max edge length
\newcommand{\Gt}{G_t}                  % heat kernel at time t
\newcommand{\Ktau}{K^{\mathrm{res}}_\tau} % resolvent (backward-Euler) kernel


% -------------------------------------------------------------------------
% Packages (include here so each document only needs % shared/macros.tex
% Shared notation for all surface-partition math documents.
% Include with: \input{../shared/macros}

% -------------------------------------------------------------------------
% Packages (include here so each document only needs \input{../shared/macros})
% -------------------------------------------------------------------------
\usepackage{amsmath,amssymb,amsthm}
\usepackage{bm}           % \bm for bold math
\usepackage{mathtools}    % \coloneqq, dcases, etc.
\usepackage{booktabs}     % \toprule etc. in tables
\usepackage{hyperref}
\usepackage{geometry}
\usepackage{microtype}
\usepackage{parskip}
\usepackage{xcolor}
\usepackage{tcolorbox}
\usepackage{enumitem}
\usepackage{float}             % [H] float placement

\geometry{a4paper, margin=2.5cm}

% -------------------------------------------------------------------------
% Theorem-like environments
% -------------------------------------------------------------------------
\theoremstyle{definition}
\newtheorem{defn}{Definition}[section]
\newtheorem{remark}{Remark}[section]

\theoremstyle{plain}
\newtheorem{prop}{Proposition}[section]

% -------------------------------------------------------------------------
% Mesh and geometry
% -------------------------------------------------------------------------
\newcommand{\V}{\mathbf{V}}            % vertex array V[i] in R^dim
\newcommand{\F}{\mathbf{F}}            % face (triangle) array
\newcommand{\vone}[1]{v_{1,#1}}       % first vertex index of VP i
\newcommand{\vtwo}[1]{v_{2,#1}}       % second vertex index of VP i

% -------------------------------------------------------------------------
% Variable points
% -------------------------------------------------------------------------
\newcommand{\lam}{\lambda}             % scalar VP parameter
\newcommand{\Lam}{\boldsymbol{\lambda}}% full parameter vector
\newcommand{\vp}[1]{\mathbf{p}_{#1}}  % 3D position of VP i
\newcommand{\ed}[1]{\mathbf{d}_{#1}}  % edge direction d_i = V[v1_i] - V[v2_i]

% -------------------------------------------------------------------------
% Steiner / triple-point
% -------------------------------------------------------------------------
\newcommand{\Sv}{\mathbf{S}}           % Steiner (Fermat-Torricelli) point
\newcommand{\Ki}{K_{i}}               % projection matrix (I-n_i n_i^T)/r_i
\newcommand{\Mmat}{M}                 % sum of K_i matrices at a triple point

% -------------------------------------------------------------------------
% Normals and unit vectors
% -------------------------------------------------------------------------
\newcommand{\nhat}{\hat{\mathbf{n}}}  % unit normal to a triangle
\newcommand{\Dhat}{\hat{\boldsymbol{\Delta}}}  % unit segment direction

% -------------------------------------------------------------------------
% Operations
% -------------------------------------------------------------------------
\newcommand{\norm}[1]{\left\|#1\right\|}
\newcommand{\abs}[1]{\left|#1\right|}
\newcommand{\ip}[2]{\langle #1,\, #2 \rangle}

% Projection perpendicular to a unit vector \uhat
\newcommand{\Proj}[1]{\bigl(\mathbf{I} - #1 #1^\top\bigr)}

% argmax operator (winner-take-all assignment in Phase 1 documents)
\DeclareMathOperator*{\argmax}{arg\,max}
\DeclareMathOperator*{\argmin}{arg\,min}

% -------------------------------------------------------------------------
% Calculus shorthands
% -------------------------------------------------------------------------
\newcommand{\pd}[2]{\dfrac{\partial #1}{\partial #2}}
\newcommand{\pdd}[3]{\dfrac{\partial^2 #1}{\partial #2\,\partial #3}}
\newcommand{\grad}{\nabla}
\newcommand{\Hess}[1]{\nabla^2 #1}

% -------------------------------------------------------------------------
% Optimization problem objects
% -------------------------------------------------------------------------
\newcommand{\Preg}{P_{\mathrm{reg}}}  % regular perimeter
\newcommand{\Pst}{P_{\mathrm{st}}}   % Steiner perimeter
\newcommand{\Areg}{A^{\mathrm{reg}}} % regular cell area
\newcommand{\Ast}{A^{\mathrm{st}}}   % Steiner area contribution
\newcommand{\Atgt}{\bar{A}}           % target area per cell
\newcommand{\Lagr}{\mathcal{L}}       % Lagrangian
\newcommand{\objfac}{\sigma}          % IPOPT objective scaling factor

% -------------------------------------------------------------------------
% Code cross-references (for "In the code..." callouts)
% -------------------------------------------------------------------------
\newcommand{\code}[1]{\texttt{#1}}
\newcommand{\codefile}[1]{\texttt{\small #1}}

% -------------------------------------------------------------------------
% Threshold dynamics / approach B (10-mbo-auction-dynamics)
% -------------------------------------------------------------------------
\newcommand{\Dlump}{D}                 % lumped mass matrix diag(v)
\newcommand{\Atau}{A_\tau}             % discrete diffusion operator (M + tau K)^{-1} M
\newcommand{\Etau}{E_\tau}             % heat-content (threshold-dynamics) energy
\newcommand{\Rcell}{R_{\mathrm{cell}}} % radius of the equal-area disc sqrt(Abar/pi)
\newcommand{\hmean}{h_{\mathrm{mean}}} % mean edge length
\newcommand{\hmax}{h_{\mathrm{max}}}   % max edge length
\newcommand{\Gt}{G_t}                  % heat kernel at time t
\newcommand{\Ktau}{K^{\mathrm{res}}_\tau} % resolvent (backward-Euler) kernel
)
% -------------------------------------------------------------------------
\usepackage{amsmath,amssymb,amsthm}
\usepackage{bm}           % \bm for bold math
\usepackage{mathtools}    % \coloneqq, dcases, etc.
\usepackage{booktabs}     % \toprule etc. in tables
\usepackage{hyperref}
\usepackage{geometry}
\usepackage{microtype}
\usepackage{parskip}
\usepackage{xcolor}
\usepackage{tcolorbox}
\usepackage{enumitem}
\usepackage{float}             % [H] float placement

\geometry{a4paper, margin=2.5cm}

% -------------------------------------------------------------------------
% Theorem-like environments
% -------------------------------------------------------------------------
\theoremstyle{definition}
\newtheorem{defn}{Definition}[section]
\newtheorem{remark}{Remark}[section]

\theoremstyle{plain}
\newtheorem{prop}{Proposition}[section]

% -------------------------------------------------------------------------
% Mesh and geometry
% -------------------------------------------------------------------------
\newcommand{\V}{\mathbf{V}}            % vertex array V[i] in R^dim
\newcommand{\F}{\mathbf{F}}            % face (triangle) array
\newcommand{\vone}[1]{v_{1,#1}}       % first vertex index of VP i
\newcommand{\vtwo}[1]{v_{2,#1}}       % second vertex index of VP i

% -------------------------------------------------------------------------
% Variable points
% -------------------------------------------------------------------------
\newcommand{\lam}{\lambda}             % scalar VP parameter
\newcommand{\Lam}{\boldsymbol{\lambda}}% full parameter vector
\newcommand{\vp}[1]{\mathbf{p}_{#1}}  % 3D position of VP i
\newcommand{\ed}[1]{\mathbf{d}_{#1}}  % edge direction d_i = V[v1_i] - V[v2_i]

% -------------------------------------------------------------------------
% Steiner / triple-point
% -------------------------------------------------------------------------
\newcommand{\Sv}{\mathbf{S}}           % Steiner (Fermat-Torricelli) point
\newcommand{\Ki}{K_{i}}               % projection matrix (I-n_i n_i^T)/r_i
\newcommand{\Mmat}{M}                 % sum of K_i matrices at a triple point

% -------------------------------------------------------------------------
% Normals and unit vectors
% -------------------------------------------------------------------------
\newcommand{\nhat}{\hat{\mathbf{n}}}  % unit normal to a triangle
\newcommand{\Dhat}{\hat{\boldsymbol{\Delta}}}  % unit segment direction

% -------------------------------------------------------------------------
% Operations
% -------------------------------------------------------------------------
\newcommand{\norm}[1]{\left\|#1\right\|}
\newcommand{\abs}[1]{\left|#1\right|}
\newcommand{\ip}[2]{\langle #1,\, #2 \rangle}

% Projection perpendicular to a unit vector \uhat
\newcommand{\Proj}[1]{\bigl(\mathbf{I} - #1 #1^\top\bigr)}

% argmax operator (winner-take-all assignment in Phase 1 documents)
\DeclareMathOperator*{\argmax}{arg\,max}
\DeclareMathOperator*{\argmin}{arg\,min}

% -------------------------------------------------------------------------
% Calculus shorthands
% -------------------------------------------------------------------------
\newcommand{\pd}[2]{\dfrac{\partial #1}{\partial #2}}
\newcommand{\pdd}[3]{\dfrac{\partial^2 #1}{\partial #2\,\partial #3}}
\newcommand{\grad}{\nabla}
\newcommand{\Hess}[1]{\nabla^2 #1}

% -------------------------------------------------------------------------
% Optimization problem objects
% -------------------------------------------------------------------------
\newcommand{\Preg}{P_{\mathrm{reg}}}  % regular perimeter
\newcommand{\Pst}{P_{\mathrm{st}}}   % Steiner perimeter
\newcommand{\Areg}{A^{\mathrm{reg}}} % regular cell area
\newcommand{\Ast}{A^{\mathrm{st}}}   % Steiner area contribution
\newcommand{\Atgt}{\bar{A}}           % target area per cell
\newcommand{\Lagr}{\mathcal{L}}       % Lagrangian
\newcommand{\objfac}{\sigma}          % IPOPT objective scaling factor

% -------------------------------------------------------------------------
% Code cross-references (for "In the code..." callouts)
% -------------------------------------------------------------------------
\newcommand{\code}[1]{\texttt{#1}}
\newcommand{\codefile}[1]{\texttt{\small #1}}

% -------------------------------------------------------------------------
% Threshold dynamics / approach B (10-mbo-auction-dynamics)
% -------------------------------------------------------------------------
\newcommand{\Dlump}{D}                 % lumped mass matrix diag(v)
\newcommand{\Atau}{A_\tau}             % discrete diffusion operator (M + tau K)^{-1} M
\newcommand{\Etau}{E_\tau}             % heat-content (threshold-dynamics) energy
\newcommand{\Rcell}{R_{\mathrm{cell}}} % radius of the equal-area disc sqrt(Abar/pi)
\newcommand{\hmean}{h_{\mathrm{mean}}} % mean edge length
\newcommand{\hmax}{h_{\mathrm{max}}}   % max edge length
\newcommand{\Gt}{G_t}                  % heat kernel at time t
\newcommand{\Ktau}{K^{\mathrm{res}}_\tau} % resolvent (backward-Euler) kernel
)
% -------------------------------------------------------------------------
\usepackage{amsmath,amssymb,amsthm}
\usepackage{bm}           % \bm for bold math
\usepackage{mathtools}    % \coloneqq, dcases, etc.
\usepackage{booktabs}     % \toprule etc. in tables
\usepackage{hyperref}
\usepackage{geometry}
\usepackage{microtype}
\usepackage{parskip}
\usepackage{xcolor}
\usepackage{tcolorbox}
\usepackage{enumitem}
\usepackage{float}             % [H] float placement

\geometry{a4paper, margin=2.5cm}

% -------------------------------------------------------------------------
% Theorem-like environments
% -------------------------------------------------------------------------
\theoremstyle{definition}
\newtheorem{defn}{Definition}[section]
\newtheorem{remark}{Remark}[section]

\theoremstyle{plain}
\newtheorem{prop}{Proposition}[section]

% -------------------------------------------------------------------------
% Mesh and geometry
% -------------------------------------------------------------------------
\newcommand{\V}{\mathbf{V}}            % vertex array V[i] in R^dim
\newcommand{\F}{\mathbf{F}}            % face (triangle) array
\newcommand{\vone}[1]{v_{1,#1}}       % first vertex index of VP i
\newcommand{\vtwo}[1]{v_{2,#1}}       % second vertex index of VP i

% -------------------------------------------------------------------------
% Variable points
% -------------------------------------------------------------------------
\newcommand{\lam}{\lambda}             % scalar VP parameter
\newcommand{\Lam}{\boldsymbol{\lambda}}% full parameter vector
\newcommand{\vp}[1]{\mathbf{p}_{#1}}  % 3D position of VP i
\newcommand{\ed}[1]{\mathbf{d}_{#1}}  % edge direction d_i = V[v1_i] - V[v2_i]

% -------------------------------------------------------------------------
% Steiner / triple-point
% -------------------------------------------------------------------------
\newcommand{\Sv}{\mathbf{S}}           % Steiner (Fermat-Torricelli) point
\newcommand{\Ki}{K_{i}}               % projection matrix (I-n_i n_i^T)/r_i
\newcommand{\Mmat}{M}                 % sum of K_i matrices at a triple point

% -------------------------------------------------------------------------
% Normals and unit vectors
% -------------------------------------------------------------------------
\newcommand{\nhat}{\hat{\mathbf{n}}}  % unit normal to a triangle
\newcommand{\Dhat}{\hat{\boldsymbol{\Delta}}}  % unit segment direction

% -------------------------------------------------------------------------
% Operations
% -------------------------------------------------------------------------
\newcommand{\norm}[1]{\left\|#1\right\|}
\newcommand{\abs}[1]{\left|#1\right|}
\newcommand{\ip}[2]{\langle #1,\, #2 \rangle}

% Projection perpendicular to a unit vector \uhat
\newcommand{\Proj}[1]{\bigl(\mathbf{I} - #1 #1^\top\bigr)}

% argmax operator (winner-take-all assignment in Phase 1 documents)
\DeclareMathOperator*{\argmax}{arg\,max}
\DeclareMathOperator*{\argmin}{arg\,min}

% -------------------------------------------------------------------------
% Calculus shorthands
% -------------------------------------------------------------------------
\newcommand{\pd}[2]{\dfrac{\partial #1}{\partial #2}}
\newcommand{\pdd}[3]{\dfrac{\partial^2 #1}{\partial #2\,\partial #3}}
\newcommand{\grad}{\nabla}
\newcommand{\Hess}[1]{\nabla^2 #1}

% -------------------------------------------------------------------------
% Optimization problem objects
% -------------------------------------------------------------------------
\newcommand{\Preg}{P_{\mathrm{reg}}}  % regular perimeter
\newcommand{\Pst}{P_{\mathrm{st}}}   % Steiner perimeter
\newcommand{\Areg}{A^{\mathrm{reg}}} % regular cell area
\newcommand{\Ast}{A^{\mathrm{st}}}   % Steiner area contribution
\newcommand{\Atgt}{\bar{A}}           % target area per cell
\newcommand{\Lagr}{\mathcal{L}}       % Lagrangian
\newcommand{\objfac}{\sigma}          % IPOPT objective scaling factor

% -------------------------------------------------------------------------
% Code cross-references (for "In the code..." callouts)
% -------------------------------------------------------------------------
\newcommand{\code}[1]{\texttt{#1}}
\newcommand{\codefile}[1]{\texttt{\small #1}}

% -------------------------------------------------------------------------
% Threshold dynamics / approach B (10-mbo-auction-dynamics)
% -------------------------------------------------------------------------
\newcommand{\Dlump}{D}                 % lumped mass matrix diag(v)
\newcommand{\Atau}{A_\tau}             % discrete diffusion operator (M + tau K)^{-1} M
\newcommand{\Etau}{E_\tau}             % heat-content (threshold-dynamics) energy
\newcommand{\Rcell}{R_{\mathrm{cell}}} % radius of the equal-area disc sqrt(Abar/pi)
\newcommand{\hmean}{h_{\mathrm{mean}}} % mean edge length
\newcommand{\hmax}{h_{\mathrm{max}}}   % max edge length
\newcommand{\Gt}{G_t}                  % heat kernel at time t
\newcommand{\Ktau}{K^{\mathrm{res}}_\tau} % resolvent (backward-Euler) kernel


\hypersetup{
  colorlinks = true,
  linkcolor  = blue!60!black,
  citecolor  = green!50!black,
  urlcolor   = blue!60!black,
  pdftitle   = {The Balanced Readout},
  pdfauthor  = {surface-partition project}
}

\title{\textbf{The Balanced Readout}\\[0.5em]
  \large Why the winner-take-all territory diverges from the constrained mass,
  and the dual-offset correction that removes the divergence}
\author{}
\date{August 2026}

\begin{document}
\maketitle

\begin{abstract}
Phase~1 minimises a $\Gamma$-convergence energy over continuous densities and
enforces equal area as a constraint on the \emph{continuous masses}. The partition
it delivers, however, is the \emph{winner-take-all territory} of those densities.
This document derives the exact identity relating the two, shows why nothing in
the variational problem controls their difference, and analyses the two artifacts
that follow --- territory far from the equal-area target, and territory that
splits into disconnected pieces. It then derives the correction implemented in
\codefile{src/partition/balanced\_readout.py}: a per-cell additive offset obtained
from the dual of a transportation problem. Three properties are proved: the
corrected readout grows each cell \emph{monotonically outward from its own core}
(the locality that in-flow corrections lacked); the attainable balance is bounded
below by the one-vertex granularity, and this bound is a theorem rather than an
observation, because the underlying linear program is \emph{not} integral; and the
offsets control area but not connectivity, which is why a second, connectivity-preserving
pass is required. Where an argument is a model rather than a derivation ---
notably the worst-cell scaling in $N$, and the mechanism behind disconnection ---
this is stated explicitly.
\end{abstract}

\tableofcontents
\newpage

\section{Overview}

The method of Bogosel and Oudet \cite{bogosel2023partitions} represents a
partition of a surface into $N$ cells by $N$ density functions and minimises a
Modica--Mortola-type functional \cite{modica1977esempio, modica1987gradient} whose
$\Gamma$-limit is the total interface length. Equal areas are imposed as
constraints on the integrals of those densities.

To hand a partition downstream one must convert the densities into a definite
assignment of each mesh vertex to one cell. The natural choice, and the one used
throughout this project, is pointwise \emph{winner-take-all}: each vertex joins
whichever cell has the largest density there. We refer to this conversion step as
the \emph{readout}, by analogy with instrumentation, where the internal state of a
device and the value it reports are distinct objects and the conversion between
them is a separate source of error.

The central observation of this document is that the constrained quantity and the
delivered quantity are \emph{different functionals of the same field}, and that
their difference is exactly computable, is supported entirely on the diffuse
interface band, and is referenced nowhere in the energy or the constraints. It is
therefore free to drift, and it drifts further as the number of cells grows at
fixed mesh resolution.

\begin{remark}
This is a statement about the readout at finite $\varepsilon$ and finite mesh. The
$\Gamma$-convergence result concerns the limit of the \emph{energy} and is not in
question. Nothing here contradicts \cite{bogosel2023partitions}; the readout is a
practical choice whose error term had not been quantified.
\end{remark}

Related documents: \codefile{docs/math/06-phase1-energy-discretization/} derives
the discrete energy itself; \codefile{docs/math/07-phase1-wta-balance/} derives a
\emph{rejected} attempt to enforce balance inside the flow, and its Proposition~1
is the identity restated as \eqref{eq:gap} below;
\codefile{docs/reference/winner\_take\_all\_partition\_gap.md} is the standing
prose account.

\section{Notation and Setup}

Let the surface be triangulated with vertex set $\{1,\dots,V\}$ and let
$\Mmat$ be the P1 finite-element mass matrix. Write
\begin{equation}
  v_i \;=\; \sum_{j} \Mmat_{ij}
  \label{eq:lumped}
\end{equation}
for the lumped mass of vertex $i$, so that $v_i > 0$ and
$\sum_i v_i = |\Sigma|$, the total surface area. The densities are
$u_{ik} \in [0,1]$ for vertex $i$ and cell $k \in \{1,\dots,N\}$, subject to the
partition-of-unity and equal-area constraints
\begin{equation}
  \sum_{k=1}^{N} u_{ik} \;=\; 1 \quad \forall i,
  \qquad
  \sum_{i=1}^{V} v_i\, u_{ik} \;=\; \Atgt \;=\; \frac{|\Sigma|}{N}
  \quad \forall k .
  \label{eq:constraints}
\end{equation}

\begin{defn}[Readout and territory]
The \emph{readout} is the map $\omega : \{1,\dots,V\} \to \{1,\dots,N\}$,
\begin{equation}
  \omega(i) \;=\; \argmax_k \, u_{ik},
  \label{eq:wta}
\end{equation}
ties broken by lowest index. The \emph{territory} of cell $k$ is
$\Omega_k = \{ i : \omega(i) = k \}$ and its \emph{discrete area} is
\begin{equation}
  T_k \;=\; \sum_{i \in \Omega_k} v_i .
  \label{eq:territory}
\end{equation}
\end{defn}

The constraint in \eqref{eq:constraints} fixes $\sum_i v_i u_{ik}$; validity is
graded on $T_k$. These coincide only if every $u_{ik}$ is exactly $0$ or $1$.

\section{The readout gap is exact}

\begin{prop}[Territory--mass identity]
For every cell $k$, with the constraint \eqref{eq:constraints} satisfied,
\begin{equation}
  T_k - \Atgt \;=\; \underbrace{\sum_{i \in \Omega_k} v_i\,(1 - u_{ik})}_{\textstyle \mathrm{gain}_k}
  \;-\; \underbrace{\sum_{i \notin \Omega_k} v_i\, u_{ik}}_{\textstyle \mathrm{lost}_k} .
  \label{eq:gap}
\end{equation}
\end{prop}

\begin{proof}
Split the constrained sum over the territory and its complement:
\[
  \Atgt \;=\; \sum_i v_i u_{ik}
        \;=\; \sum_{i \in \Omega_k} v_i u_{ik} \;+\; \sum_{i \notin \Omega_k} v_i u_{ik}
        \;=\; \Big( T_k - \sum_{i \in \Omega_k} v_i (1-u_{ik}) \Big) \;+\; \mathrm{lost}_k ,
\]
using $\sum_{i \in \Omega_k} v_i u_{ik} = T_k - \sum_{i \in \Omega_k} v_i(1-u_{ik})$
from \eqref{eq:territory}. Rearranging gives \eqref{eq:gap}.
\end{proof}

The identity is exact: no linearisation, no assumption on $\varepsilon$ or on the mesh.
Its two terms have a direct reading. $\mathrm{gain}_k$ is mass that cell $k$ \emph{owns
but does not hold} --- vertices it wins while its density there is below one.
$\mathrm{lost}_k$ is mass it \emph{holds but does not own} --- density deposited at
vertices won by somebody else.

\begin{remark}[The gap lives entirely in the interface band]
Both terms vanish wherever the field is crisp. Deep inside cell $k$,
$u_{ik} \approx 1$, so $1 - u_{ik} \approx 0$ and the vertex contributes nothing
to $\mathrm{gain}_k$. Far outside, $u_{ik} \approx 0$ and it contributes nothing to
$\mathrm{lost}_k$. Only vertices in the diffuse band, where $0 < u_{ik} < 1$, contribute
to either. \emph{The readout gap is a band phenomenon}, and everything about how
it scales follows from how much of a cell the band occupies.
\end{remark}

\subsection{Why nothing controls it}

The energy is a function of $u$; the constraints \eqref{eq:constraints} are linear
in $u$. Neither references $\omega$. Moreover the map $u \mapsto \omega$ is
piecewise constant and therefore discontinuous: it is constant on each cell of the
arrangement $\{u_{ik} - u_{il} = 0\}$ and jumps across it. Consequently
\begin{itemize}\itemsep2pt
  \item $T_k$ has zero gradient with respect to $u$ almost everywhere, so no
    first-order optimality condition can constrain it;
  \item adding $T_k$ to the objective is not possible without a surrogate, and the
    surrogate used in \codefile{docs/math/07-phase1-wta-balance/} was measured and
    rejected (see \S\ref{sec:locality});
  \item the quantity $\mathrm{gain}_k - \mathrm{lost}_k$ is, from the optimiser's point of view,
    an unobserved by-product.
\end{itemize}

\section{Artifact I: territory far from the equal-area target}
\label{sec:artifact1}

\subsection{The band fraction}

Let $R_{\mathrm{cell}} = \sqrt{|\Sigma| / (\pi N)}$ be the radius of a disc of the
target area and let $h$ be a representative edge length. In the discretisation
used here the interface width parameter is tied to the mesh, $\varepsilon \sim h$
(\codefile{docs/math/06-phase1-energy-discretization/}). A cell's band is a collar
of width $O(\varepsilon)$ around a boundary of length $O(R_{\mathrm{cell}})$, so the
fraction of a cell occupied by band is
\begin{equation}
  f_b \;\sim\; \frac{\varepsilon \cdot 2\pi R_{\mathrm{cell}}}{\pi R_{\mathrm{cell}}^2}
     \;=\; \frac{2\varepsilon}{R_{\mathrm{cell}}}
     \;\sim\; \frac{h}{R_{\mathrm{cell}}} .
  \label{eq:bandfrac}
\end{equation}
With $h \sim \sqrt{|\Sigma|/V}$ this becomes
\begin{equation}
  f_b \;\sim\; \sqrt{\frac{N}{V}} \;=\; \Big(\frac{V}{N}\Big)^{-1/2},
  \label{eq:bandscaling}
\end{equation}
i.e.\ \textbf{the band fraction is governed by vertices per cell alone.} This
recovers, from geometry, the empirical form $f_b = 3.72 / \sqrt{2V/N}$ recorded in
\codefile{docs/reference/winner\_take\_all\_partition\_gap.md}.

\begin{remark}[The scaling trap]
At \emph{fixed mesh}, raising $N$ shrinks $R_{\mathrm{cell}}$ and therefore
\emph{raises} $f_b$: the band becomes a larger share of each cell precisely as the
cells become the objects one cares about. Holding $f_b$ fixed requires
$V \propto N$ --- constant vertices per cell --- so the mesh must grow linearly in
the number of cells merely to keep the readout error from worsening.
\end{remark}

\subsection{The diffuse runt}

Equation \eqref{eq:gap} permits a configuration that is feasible, low-energy, and
useless: a cell whose density is spread as a broad halo, everywhere slightly below
its neighbours'. It then satisfies $\sum_i v_i u_{ik} = \Atgt$ exactly while
winning almost nothing, because $\mathrm{lost}_k$ is large and $\mathrm{gain}_k$ small. Such a
cell is \emph{confidently in the right place and insufficiently loud}. Empirically
at $N=300$ the worst cell reached $36.15\%$ below target while its continuous mass
was exact.

\subsection{Worst-of-$N$ --- a model, not a derivation}

\begin{remark}[Explicitly not established]
The following is the argument by which the project has explained why the problem
sharpens with $N$. It is a \textbf{model}. The distribution of $\mathrm{gain}_k - \mathrm{lost}_k$
across cells has never been measured, and the available data points differ in both
$\lambda$ and mesh, so they cannot test it.
\end{remark}

If one models each cell's net exchange as a random variable with mean zero and
standard deviation $\sigma \propto f_b$, roughly independent across cells, then
validity --- which is graded on $\max_k |T_k - \Atgt|$ --- is an extreme-value
statistic, and for light-tailed exchange the maximum of $N$ draws grows like
$\sigma\sqrt{2\ln N}$. Two consequences follow \emph{if the model holds}: the worst
cell degrades slowly but without bound as $N$ grows even at fixed $f_b$; and any
mechanism that merely \emph{reduces the variance} of the exchange postpones the
problem rather than removing it. Only a mechanism that makes the delivered
representation balanced \emph{by construction} is immune to the tail. That
prediction is consistent with every measured outcome in this project, which is
evidence for the model but not a proof of it.

\section{Artifact II: disconnected territory}
\label{sec:artifact2}

\begin{remark}[Weaker than \S\ref{sec:artifact1}, and deliberately so]
The existence of this artifact is a theorem (\S\ref{sec:disc-possible}); its
\emph{cause} in our runs is a hypothesis that has never been traced to a specific
event. The two are separated below.
\end{remark}

\subsection{Nothing forbids it}
\label{sec:disc-possible}

The readout \eqref{eq:wta} is defined pointwise. $\Omega_k$ is a super-level set
of the function $i \mapsto u_{ik} - \max_{l \neq k} u_{il}$, and a super-level set
of an arbitrary function on a graph need not be connected. Neither the energy nor
the constraints contain any term that is finite for connected territories and
infinite otherwise. Disconnection is therefore \emph{admissible} and costs nothing
directly; a split cell is a legitimate local minimiser of the relaxed problem.

What makes it pathological is external: a minimal-perimeter cell is connected, so
a split territory is a non-physical artifact of the representation, and Phase~2
cannot consume it, since a split cell yields multi-loop contours.

\begin{remark}[Why the other two gates are blind to it]
A split cell's pieces sum to the correct area and each piece is crisp. It
therefore passes a dormant-cell check (peak density $\approx 1$) and an
area-imbalance check (total area correct). Connectivity is a third, independent
validity dimension.
\end{remark}

\subsection{The mechanism --- hypothesis}

The constraint $\sum_i v_i u_{ik} = \Atgt$ is enforced by an iterative projection
onto the feasible set. The projection is a \emph{nonlocal} operator: correcting a
deficit in cell $k$ adds a rank-one term spread across the whole free set,
followed by renormalisation. Raising $u_{ik}$ everywhere by a near-uniform amount
can therefore flip the argmax at vertices arbitrarily far from the cell's core ---
wherever cell $k$ happened to be a close second --- and open a second basin.

Three observations are consistent with this and none establishes it:
\begin{enumerate}\itemsep2pt
  \item Splits appear \emph{while balance is improving}. At $N=300$ they are born
    mid-level-2, at the same time as the imbalanced count falls from 234 to 10:
    balance is being bought with disconnection.
  \item Mechanisms that push balance through the projection make it dramatically
    worse. A discrete-area trim produced 14/200 split cells against a matched
    control's 0/200; a soft area constraint produced 3--4/100.
  \item The correction derived in \S\ref{sec:correction}, which is
    \emph{local} by Proposition~\ref{prop:nested}, does not manufacture splits.
\end{enumerate}

\begin{remark}[What would settle it]
Decompose the per-iteration update into its gradient and projection contributions
at the iteration a split is born, and test whether the new component's vertices had
$u_{ik}$ raised by the projection. This has not been done; it is recorded as future
work in \codefile{docs/plans/PUBLICATION\_READINESS\_PLAN.md}.
\end{remark}

\section{The corrected readout}
\label{sec:correction}

\subsection{The assignment problem}

Fix any per-vertex score matrix $s_{ik}$; the shipped readout uses
$s_{ik} = \log u_{ik}$, with densities clipped from below so that exact zeros give
a finite, very negative score. We seek a hard assignment maximising total score
subject to every cell attaining its target area:
\begin{equation}
  \max_{x} \; \sum_{i,k} v_i\, s_{ik}\, x_{ik}
  \quad\text{s.t.}\quad
  \sum_k x_{ik} = 1 \;\; \forall i, \qquad
  \sum_i v_i\, x_{ik} = \Atgt \;\; \forall k, \qquad
  x_{ik} \ge 0 .
  \label{eq:lp}
\end{equation}
This is a transportation problem: supplies $v_i$ at vertices, demands $\Atgt$ at
cells, and both marginals sum to $|\Sigma|$, so it is feasible.

\begin{remark}[Terminology]
This project's prose has called \eqref{eq:lp} a \emph{semi-discrete} optimal
transport problem. Strictly it is \emph{discrete-to-discrete}: $V$ weighted points
to $N$ cells. The semi-discrete problem --- a continuous source measure
transported to $N$ atoms, in the sense of
\cite{aurenhammer1998minkowski, merigot2011multiscale, kitagawa2019convergence} ---
is its continuum limit as the mesh refines. The offsets derived below are the
discrete analogue of the power weights in that literature.
\end{remark}

\subsection{The dual, and where the offsets come from}

Introduce multipliers $\varphi_i$ for the row constraints and $\psi_k$ for the
column constraints. Stationarity of the Lagrangian in $x_{ik}$ gives
$v_i s_{ik} - \varphi_i - v_i \psi_k \le 0$, with equality wherever $x_{ik} > 0$,
so $\varphi_i = v_i \max_k (s_{ik} - \psi_k)$ and the dual is
\begin{equation}
  \min_{\psi \in \mathbb{R}^N} \; D(\psi), \qquad
  D(\psi) \;=\; \sum_i v_i \max_k \big( s_{ik} - \psi_k \big) \;+\; \Atgt \sum_k \psi_k .
  \label{eq:dual}
\end{equation}
$D$ is a sum of pointwise maxima of affine functions plus a linear term, hence
\textbf{convex} and piecewise linear. Complementary slackness forces every vertex
onto a maximising cell, so an optimal assignment is
\begin{equation}
  \omega_\psi(i) \;=\; \argmax_k \big( s_{ik} - \psi_k \big) .
  \label{eq:shifted}
\end{equation}

\begin{remark}[Sign convention]
The implementation stores $\psi^{\mathrm{code}} = -\psi$ and evaluates
$\argmax_k ( s_{ik} + \psi^{\mathrm{code}}_k )$. All statements below use the code
convention, in which \emph{raising} $\psi_k$ \emph{grows} cell $k$.
\end{remark}

\begin{prop}[The implemented update is the dual subgradient]
\label{prop:subgrad}
At any $\psi$ where the argmax is unique at every vertex, $D$ is differentiable
with
\begin{equation}
  \frac{\partial D}{\partial \psi_k} \;=\; \Atgt - T_k(\psi),
  \qquad T_k(\psi) = \sum_{i \,:\, \omega_\psi(i) = k} v_i ,
  \label{eq:subgrad}
\end{equation}
and in general this is a subgradient. Hence gradient descent on $D$ is exactly
``grow the cells that are too small, shrink the cells that are too large''.
\end{prop}

\begin{proof}
Differentiating the first term of \eqref{eq:dual} at a point of uniqueness picks
out, for each $i$, the maximising index $\omega_\psi(i)$, contributing $-v_i$ to
the $\psi_k$ derivative precisely when $\omega_\psi(i) = k$; the second term
contributes $\Atgt$. Summing gives \eqref{eq:subgrad}. At points where the argmax
is non-unique, $D$ is a maximum of finitely many affine functions and the same
expression is an element of the subdifferential for any measurable tie-breaking.
\end{proof}

\begin{tcolorbox}[colback=blue!4!white, colframe=blue!40!black,
                  title={Code: dual ascent on the offsets}, fontupper=\small]
\codefile{src/partition/balanced\_readout.py}: \code{solve\_dual\_offsets(scores, lumped\_mass, target, config)}\\[2pt]
\code{psi -= (eta0 / (1 + decay*it)) * rel}, with
\code{rel = (areas - target)/target} --- the normalised form of
\eqref{eq:subgrad}. The routine returns its \emph{best} iterate, since $D$ is
piecewise linear and subgradient descent is not monotone. Scores enter only
through \code{argmax(scores + psi)}, so the routine is correct for any additive
score matrix; only its step schedule is scale-dependent.
\end{tcolorbox}

\subsection{Locality: cells grow outward from their own cores}
\label{sec:locality}

\begin{prop}[Nested monotone growth]
\label{prop:nested}
Fix $\psi_{-k}$ and regard $\Omega_k(\psi_k)$ as a function of the single offset
$\psi_k$. Define the \emph{margin} of vertex $i$ against cell $k$ as
\begin{equation}
  m_i \;=\; \max_{l \neq k} \big( s_{il} + \psi_l \big) \;-\; s_{ik} .
  \label{eq:margin}
\end{equation}
Then $\Omega_k(\psi_k) = \{ i : m_i \le \psi_k \}$. Consequently
$\psi_k < \psi_k'$ implies $\Omega_k(\psi_k) \subseteq \Omega_k(\psi_k')$, and
vertices are acquired in nondecreasing order of $m_i$.
\end{prop}

\begin{proof}
Immediate from \eqref{eq:shifted}: $i \in \Omega_k$ iff
$s_{ik} + \psi_k \ge \max_{l\neq k}(s_{il} + \psi_l)$, i.e.\ iff $\psi_k \ge m_i$.
The sets are sublevel sets of the fixed function $m$, hence nested.
\end{proof}

This is the property the in-flow corrections lacked, and it is worth stating
plainly what it does and does not say.

\begin{itemize}\itemsep3pt
  \item \textbf{What it gives.} Growth is \emph{ordered by margin}. A cell claims
    first the vertices where it was most nearly winning, which --- for a density
    field that is unimodal about the cell's core --- are the vertices adjacent to
    the territory it already holds. Correction proceeds outward through the
    adjacent band rather than being deposited wherever the algebra puts it.
  \item \textbf{What it does not give.} Proposition~\ref{prop:nested} is a
    statement about the \emph{order} of acquisition, not about connectivity.
    If the margin function $m$ has two separated basins --- exactly the situation
    Artifact~II describes --- then nested growth will populate both. Locality
    prevents the correction from \emph{manufacturing} new components; it cannot
    repair components that the density field already possesses. Hence
    \S\ref{sec:repair}.
\end{itemize}

Contrast this with enforcing balance inside the flow. There the correction reaches
the field through the constraint projection, which is nonlocal
(\S\ref{sec:artifact2}); measured outcomes were a net regression in both attempts
(\codefile{docs/experiments/04-territory-aware-highn-validation/},
\codefile{docs/experiments/05-soft-area-constraint/}).

\subsection{The granularity floor is a theorem, not an observation}

It is tempting to assert that \eqref{eq:lp} has an integral optimum because
transportation problems do. That reasoning is unavailable here.

\begin{prop}[Non-integrality]
\label{prop:integrality}
The classical integrality theorem for the transportation polytope requires
integer supplies and demands. In \eqref{eq:lp} the supplies are the real lumped
masses $v_i$ and the demands are $\Atgt = |\Sigma|/N$. Unless some subset
$S \subseteq \{1,\dots,V\}$ satisfies $\sum_{i \in S} v_i = \Atgt$ exactly, no
integral feasible point exists at all, and every vertex of the polytope splits at
least one vertex's mass between cells. A basic feasible solution of a
transportation problem with $V$ sources and $N$ sinks has at most $V + N - 1$
nonzeros, so at most $N-1$ vertices are split.
\end{prop}

Consequently the exact optimum of \eqref{eq:lp} is in general \emph{fractional},
and the readout \eqref{eq:shifted} is a rounding of it. The residue is bounded by
the mass of the split vertices, which yields the floor
\begin{equation}
  \min_{\psi} \; \max_k \frac{\big| T_k(\psi) - \Atgt \big|}{\Atgt}
  \;\ge\; \text{(generically)} \;\; \frac{\max_i v_i}{\Atgt}
  \;\approx\; \frac{N}{V} \cdot \frac{\max_i v_i}{\bar v} ,
  \label{eq:granularity}
\end{equation}
where $\bar v = |\Sigma|/V$. \textbf{A vertex is indivisible, so no assignment
method whatever can do better.} This is why acceptance thresholds in this project
are stated as multiples of $\max_i v_i / \Atgt$ rather than as absolute
percentages: a bar below \eqref{eq:granularity} is not demanding, it is
impossible.

\begin{remark}[Measured values]
$\max_i v_i / \Atgt$ is $1.011\%$ at $V=47{,}488$, $N=300$; $0.421\%$ at
$V=114{,}144$, $N=300$; $0.140\%$ at $V=114{,}144$, $N=100$; and $0.561\%$ at
$V=114{,}144$, $N=400$. The floor scales as $N/V$ --- inversely with vertices per
cell --- for the same reason \eqref{eq:bandscaling} does.
\end{remark}

\section{Connectivity repair}
\label{sec:repair}

The offsets control area but, by the second bullet of \S\ref{sec:locality}, not
connectivity. A second pass restores it. It operates on labels only.

\paragraph{Stage 1 --- island absorption.} For each cell, compute the connected
components of its territory in the mesh edge graph. Every non-largest component
whose area exceeds a fixed fraction of $\Atgt$ is relabelled to the neighbouring
cell with which it shares the longest boundary. This step is unconditionally safe:
a component adjacent to cell $c$ can always be merged into $c$ without
disconnecting $c$, since the merged set is connected through the shared boundary.
On exit every cell is connected (or empty, which is reported).

\paragraph{Stage 2 --- area restoration by boundary transfer.} Absorption
perturbs the areas, so equal areas are restored by moving \emph{single boundary
vertices} from richer cells to poorer ones. A move of vertex $i$ from donor $d$ to
receiver $r$ is admitted only if
\begin{equation}
  T_d - T_r \;>\; v_i ,
  \label{eq:improving}
\end{equation}
which guarantees that the transfer strictly reduces $|T_d - \Atgt| + |T_r - \Atgt|$
rather than overshooting, and only if $\Omega_d \setminus \{i\}$ remains connected
--- an exact articulation test, cheap because it is local to one cell. Candidates
are ranked by the score margin $s_{ir} - s_{id}$, so the least costly transfers
are taken first.

\begin{remark}[What is and is not guaranteed]
Both invariants --- every cell connected, and no admissible improving move
remaining --- hold on exit \emph{by construction}. What is \textbf{not}
guaranteed is that the exit state attains the granularity floor
\eqref{eq:granularity}: stage 2 is a greedy local search over single-vertex moves,
and a configuration can be locally optimal without being globally balanced,
particularly when \eqref{eq:improving} blocks every candidate. The implementation
therefore reports its own strain --- number of moves, moves blocked by the
connectivity test, sweeps used, and whether the sweep cap was reached --- so that
a run beyond what repair can fix says so rather than returning a degraded
partition silently.
\end{remark}

\begin{tcolorbox}[colback=blue!4!white, colframe=blue!40!black,
                  title={Code: the two-stage readout}, fontupper=\small]
\codefile{src/partition/balanced\_readout.py}:
\code{apply\_balanced\_readout(...)} orchestrates
\code{solve\_dual\_offsets} $\to$ \code{reassign\_islands} $\to$
\code{rebalance\_boundaries}; \code{\_stays\_connected} is the articulation test
and \code{build\_vertex\_adjacency} the edge graph, built from \code{faces} so
that it respects the true surface topology (e.g.\ the torus periodic wrap) rather
than a flat parametrisation. CLI: \codefile{scripts/balanced\_readout.py}.
\end{tcolorbox}

\section{Relation to power diagrams}

With scores $s_{ik} = -d^2(i, z_k)$ for sites $z_k$, \eqref{eq:shifted} reads
$\omega(i) = \argmax_k ( -d^2(i,z_k) + \psi_k ) = \argmin_k ( d^2(i,z_k) - \psi_k )$,
which is precisely a \emph{power diagram} (Laguerre tessellation) with weights
$\psi_k$. That the weights can be chosen to prescribe the cell capacities is
Aurenhammer, Hoffmann and Aronov's Minkowski-type theorem
\cite{aurenhammer1998minkowski}; modern solvers for the continuous version are
given by Mérigot \cite{merigot2011multiscale} and Kitagawa, Mérigot and Thibert
\cite{kitagawa2019convergence}, and the construction is used for blue-noise
sampling by de Goes \emph{et al.} \cite{degoes2012bluenoise}.

The balanced readout is the same algebra applied to a different score. Two
differences are worth recording:
\begin{enumerate}\itemsep2pt
  \item Our scores are $\log u$, a \emph{relaxation output}, not a distance. There
    are no sites, and no geometric structure is assumed --- only that the score is
    additive in $\psi$.
  \item Our problem is finite-to-finite, so the damped Newton methods of
    \cite{kitagawa2019convergence}, which exploit the smoothness of the
    semi-discrete objective in the continuous-source setting, do not transfer
    directly; $D$ in \eqref{eq:dual} is piecewise linear, hence nonsmooth.
    Subgradient descent is used instead. Entropic regularisation
    \cite{cuturi2013sinkhorn} is an alternative that was implemented and measured
    for this problem, and lost to subgradient descent on five of six score
    matrices (\codefile{docs/experiments/07-phase0-shared-harness/}).
\end{enumerate}

\begin{remark}[We are not aware of this being reported]
Applying a capacity-constrained assignment as a \emph{readout correction for a
$\Gamma$-convergence relaxation} --- as opposed to as a construction in its own
right --- is not something we have found in the literature. Stated as a population
claim, the corpus searched is this project's bibliography plus two full-text
readings; that is not a systematic review, and the claim is phrased accordingly.
\end{remark}

\section{Measured effect}

Recorded here for completeness; the study that measures them belongs in
\codefile{docs/experiments/}, which does not yet contain one for this stage.

\begin{center}\small
\begin{tabular}{@{}lrrr@{}}
\toprule
torus, $N=300$, $V=47{,}488$ & source & after readout & \\
\midrule
imbalanced cells (gate $5\%$) & 10 & \textbf{0} & \\
worst cell deviation & $36.15\%$ & $\mathbf{1.63\%}$ & floor $1.011\%$ \\
disconnected cells & 2 & \textbf{0} & \\
label-boundary length & $189.73$ & $193.30$ & $+1.88\%$ \\
wall time & --- & $\approx 17$\,s & \\
\bottomrule
\end{tabular}
\end{center}

The $+1.88\%$ boundary cost is the price of the correction: territory is moved to
meet the area target, and moved territory has a longer border. Phase~2 subsequently
recovers most of it.

\section*{Summary table}

\begin{center}\small
\begin{tabular}{@{}llll@{}}
\toprule
Quantity & Status & Where & Implementation \\
\midrule
Territory--mass identity \eqref{eq:gap} & derived, exact & \S3 & --- \\
Band fraction \eqref{eq:bandscaling} & derived, scaling & \S\ref{sec:artifact1} & --- \\
Worst-of-$N$ growth & \textbf{model, unmeasured} & \S\ref{sec:artifact1} & --- \\
Disconnection admissible & derived & \S\ref{sec:disc-possible} & --- \\
Disconnection \emph{mechanism} & \textbf{hypothesis, untraced} & \S\ref{sec:artifact2} & --- \\
Dual \eqref{eq:dual}, subgradient \eqref{eq:subgrad} & derived & \S\ref{sec:correction} & \code{solve\_dual\_offsets} \\
Nested growth (locality) & derived & \S\ref{sec:locality} & --- \\
Granularity floor \eqref{eq:granularity} & derived from non-integrality & \S\ref{sec:correction} & \code{vertex\_granularity} \\
Island absorption & derived, safe & \S\ref{sec:repair} & \code{reassign\_islands} \\
Boundary transfer & greedy; invariants only & \S\ref{sec:repair} & \code{rebalance\_boundaries} \\
\bottomrule
\end{tabular}
\end{center}

\bibliographystyle{plain}
\bibliography{../shared/references}

\end{document}
